% s02-5-multi.tex —— 2.3.8 小节：多点处的 Iwahori 层级结构（印刷页 54–55）

\subsection{多点处的 Iwahori 层级结构}
\label{sec:multi}

全文迄今考虑的是 $\Bun^I_G$，即曲线 $X$ 上带单点 $x\in X$ 处 Iwahori 层级结构的主
$G$-丛模叠。其实也可以考虑 $k$ 个互异点 $x_1,\dots,x_k\in X$ 处的 Iwahori 层级结构。
本小节说明：为证明相应模叠上的 $D$-模范畴紧生成，证明需要作哪些改动。

\begin{definition}
\label{def:multi-bunI}
设 $x_1,\dots,x_k\in X$ 为 $k$ 个互异点。定义
\[
\Bun^{(I;x_1,\dots,x_k)}_G=\Bun_G\times_{(\pt/G)^k}(\pt/B)^k,
\]
其中映射为
\[
(x_1^*,\dots,x_k^*)\colon \Bun_G\to (\pt/G)^k .
\]
\end{definition}

\begin{proposition}
\label{prop:multi-cg}
$D\text{-}\mathrm{mod}(\Bun^{(I;x_1,\dots,x_k)}_G)$ 紧生成。
\end{proposition}

可以看见，命题~\ref{prop:multi-cg} 的证明在概念上与单点情形的证明完全相同，
差别只在数字上。正如从无 Iwahori 分歧的情形（即 [DG15]）过渡到单点情形时，
正确的界从 $2g-2$ 增大到 $2g-1$，$k$ 点情形中正确的界是 $2g-2+k$。
\srcpage{54}
换言之，证明在于说明：只要 $\langle\theta,\check\alpha_i\rangle\ge 2g-2+k$，拟紧开子叠
\[
\Bun^{(I;x_1,\dots,x_k),(\le\theta)}_G
\]
就余截断。做法是把
\[
\Bun^{(I;x_1,\dots,x_k),(\le\theta')}_G
\setminus \Bun^{(I;x_1,\dots,x_k),(\le\theta)}_G
\]
写成有限个 $\Bun^{(I;x_1,\dots,x_k),(S)}_G$ 之并，其中 $S$ 是 $k$-Iwahori 收缩
$P$-可容的（$P$ 可变动），再按
\[
{}^{(w_1,\dots,w_k)}\Bun^{(I;x_1,\dots,x_k),(S)}_G
=\Bun^{(I;x_1,\dots,x_k),(S)}_G
\times_{(P\backslash G/B)^k}
\bigl(P\backslash Pw_1B/B\times\cdots\times P\backslash Pw_kB/B\bigr)
\]
进一步分层。理由归根结底是：此处需要的不是 $n(P)_{F_M}$ 的全局生成，而是它的
\emph{$k$-重全局生成}，即要求
\[
H^0(X,n(P)_{F_M})\to \bigoplus_{j=1}^{k}(n(P)_{F_M})_{x_j}
\]
对 $X$ 的任意 $k$ 个点都满。
\srcpage{55}
